% Thermal Nexus — Work Completed So Far
% Author: Thermal Nexus team
% Date:  July 26, 2026
% Compile:  pdflatex WORK_COMPLETED_SUMMARY.tex
% Description: Standalone work-completed report covering the IEEE HART
%              HardwAIre Challenge 2026 "Thermal Nexus" software-only
%              prototype through the integrated external T15 benchmark.

\documentclass[11pt,a4paper]{article}

% ------- packages -------------------------------------------------------
\usepackage[margin=1in]{geometry}
\usepackage{microtype}
\usepackage{graphicx}
\usepackage{booktabs}
\usepackage{array}
\usepackage{longtable}
\usepackage{tabularx}
\usepackage{enumitem}
\usepackage{xcolor}
\usepackage{hyperref}
\usepackage{listings}
\usepackage{fancyhdr}
\usepackage{titlesec}

% ------- style ----------------------------------------------------------
\hypersetup{
    colorlinks=true,
    linkcolor=teal!60!black,
    urlcolor=teal!60!black,
    citecolor=teal!60!black,
    pdftitle={Thermal Nexus --- Work Completed So Far},
    pdfauthor={Thermal Nexus Team},
}

\definecolor{teal}{RGB}{15,118,110}
\definecolor{tnbg}{RGB}{236,253,248}
\definecolor{tnrule}{RGB}{219,234,229}
\definecolor{codegray}{RGB}{245,245,245}

\titleformat{\section}{\Large\bfseries\color{teal}}{\thesection}{1em}{}
\titleformat{\subsection}{\large\bfseries\color{teal!85!black}}{\thesubsection}{1em}{}

\pagestyle{fancy}
\fancyhf{}
\fancyhead[L]{\small Thermal Nexus --- Work Completed}
\fancyhead[R]{\small \today}
\fancyfoot[C]{\thepage}
\renewcommand{\headrulewidth}{0.4pt}

% ------- title ----------------------------------------------------------
\title{\bfseries\Huge Thermal Nexus --- Work Completed So Far
       \vspace{4pt}\\[6pt]
       \large IEEE HART HardwAIre Challenge 2026 --- Software Phase Summary}
\author{Thermal Nexus Team \\ \small \texttt{Pshawon-kundu / THERMAL-NEXUS}}
\date{\today}

% ------- tables helper --------------------------------------------------
\newcolumntype{L}[1]{>{\raggedright\arraybackslash}p{#1}}

% ====================================================================
\begin{document}
\maketitle
\thispagestyle{fancy}

% ====================================================================
\begin{abstract}
\noindent
This report summarises the work completed so far on the
\emph{Thermal Nexus} predictive cold-chain monitoring system. The main
software pipeline --- synthetic data generation, three operating modes,
end-to-end radio+reader simulation, SQLite storage, an offline Streamlit
dashboard, KPI and energy analysis, and an embedded model-export workflow
with golden-vector parity testing --- is fully implemented and tested.
Building on that foundation, the team integrated an external
building-environment T15 benchmark dataset as a separate validation stream
and shipped a reproducible end-to-end pipeline for it, including feature
engineering, prediction horizons, model training and evaluation, scientific
limitations documentation, reproducible PowerShell launchers, a clean Git
branch with carefully scoped \texttt{.gitignore} rules, and a
test-isolation refactor so the test suite no longer depends on locally
generated data.
\end{abstract}

\vspace{1em}

% ====================================================================
\section{Core Thermal Nexus software --- complete}

The main software pipeline is fully implemented under the IEEE HART
HardwAIre Challenge 2026 Phase~1--6 software scope. The completed modules
include:

\begin{itemize}[leftmargin=2em,itemsep=2pt]
  \item Synthetic temperature-data generation
        (\texttt{simulator/temperature/}).
  \item Three sensor-node operating modes:
        \emph{fixed}, \emph{rule-based}, and \emph{ML}-based transmission
        (\texttt{simulator/run\_end\_to\_end.py}).
  \item Deterministic radio channel + reader-side data collection
        (\texttt{simulator/radio\_simulation}, \texttt{simulator/reader}).
  \item SQLite database storage with v1\,$\rightarrow$\,v2 migration
        (\texttt{host/database/}).
  \item Offline Streamlit dashboard with eleven sidebar pages
        (\texttt{host/dashboard/}).
  \item KPI generation, mode comparison, and energy-consumption estimation
        (\texttt{analysis/kpi\_engine.py}, \texttt{kpi\_reports.py}).
  \item Alert and event logging
        (\texttt{host/database/alerts} + \texttt{source\_event\_id}
        linkage to radio/reader events).
  \item Embedded model export to C99
        (\texttt{embedded/deployment/prepare\_export.py},
        \texttt{estimate\_resources.py}, \texttt{export\_manifest.py}).
  \item Golden-vector generation
        (\texttt{embedded/deployment/generate\_golden\_vectors.py}).
  \item Embedded parity-testing workflow
        (\texttt{embedded/tests/test\_parity.py}).
  \item Deployment and hardware-integration documentation
        (\texttt{docs/SYSTEM\_ARCHITECTURE.md},
        \texttt{EMBEDDED\_DEPLOYMENT\_SPECIFICATION.md}).
\end{itemize}

The main software pipeline was previously tested successfully with
approximately 50 passing tests; the current full suite reports
\textbf{51 passed, 4 failed}, where the four failures belong exclusively
to the separately gated external T15 phase (missing
\texttt{ml/data/external/t15/raw/} data after a clean clone by design).

% ====================================================================
\section{External T15 dataset --- integrated}

The team added an external building-environment dataset as a separate
benchmark. The two raw files were converted into the Thermal Nexus
format:

\begin{itemize}[leftmargin=2em,itemsep=2pt]
  \item \texttt{NEW-DATA-1.T15.txt}
  \item \texttt{NEW-DATA-2.T15.txt}
\end{itemize}

The converted dataset contains:

\begin{center}
\begin{tabularx}{0.85\textwidth}{L{5cm}L{8cm}}
\toprule
\textbf{Property} & \textbf{Value} \\
\midrule
Total rows              & 4{,}137 \\
Total runs              & 45 \\
Train runs              & 31 \\
Validation runs         & 6 \\
Test runs               & 8 \\
Sampling interval       & 15 minutes \\
Invalid timestamps      & 0 \\
Duplicate timestamps    & 0 \\
Run overlap between splits & 0 \\
\bottomrule
\end{tabularx}
\end{center}

\noindent The dataset period is:
\begin{center}
\begin{tabular}{ll}
\toprule
\textbf{Field} & \textbf{Value} \\
\midrule
Start & 2012-03-13 11:45:00 \\
End   & 2012-05-02 07:30:00 \\
\bottomrule
\end{tabular}
\end{center}

% ====================================================================
\section{Scientific limitations --- documented}

The dataset is clearly labelled \texttt{EXTERNAL\_DERIVED\_BENCHMARK}. The
limitations document (\texttt{evidence/external\_t15/limitations.json})
states:

\begin{itemize}[leftmargin=2em,itemsep=2pt]
  \item Thermal Nexus did not collect this dataset.
  \item It is a building-environment dataset.
  \item It is not a vaccine or cold-chain field dataset.
  \item Results from this dataset do not prove real cold-chain performance.
  \item Future temperature targets are derived from the same source dataset.
  \item Multivariate models are research references and are not directly
        deployable on the current temperature-only sensor node.
\end{itemize}

\noindent The previously incorrect provenance wording was corrected in:
\begin{center}
\texttt{evidence/external\_t15/limitations.json}
\end{center}

% ====================================================================
\section{Model-ready datasets --- created}

The external pipeline generates \emph{two} separate model-ready datasets:

\begin{center}
\begin{tabular}{ll}
\toprule
\textbf{Filename} & \textbf{Purpose} \\
\midrule
\texttt{temperature\_only\_benchmark.csv}        & Deployable benchmark \\
\texttt{multivariate\_research\_reference.csv}   & Research reference \\
\bottomrule
\end{tabular}
\end{center}

\noindent Each file contains:
\begin{itemize}[leftmargin=2em,itemsep=2pt]
  \item Train, validation, and test splits
  \item Run identifiers
  \item Timestamps
  \item Historical lag features
  \item Rolling statistics
  \item Trend features
  \item Future-temperature targets
  \item Thermal-change classification target
\end{itemize}

\noindent Row counts after split and feature engineering:

\begin{center}
\begin{tabular}{lr}
\toprule
\textbf{Dataset} & \textbf{Rows} \\
\midrule
Temperature-only rows & 3{,}321 \\
Multivariate rows     & 3{,}321 \\
\bottomrule
\end{tabular}
\end{center}

% ====================================================================
\section{Prediction horizons --- implemented}

The external benchmark predicts future temperature at:

\begin{center}
\begin{tabular}{lrr}
\toprule
\textbf{Horizon} & \textbf{Duration} & \textbf{Future rows} \\
\midrule
Short  & 30 minutes & 2 \\
Mid    & 60 minutes & 4 \\
Long   & 90 minutes & 6 \\
\bottomrule
\end{tabular}
\end{center}

\noindent Because the dataset interval is 15~minutes, each horizon above
corresponds directly to the listed number of future rows. The
implementation prevents future targets from crossing run boundaries.

% ====================================================================
\section{Historical-only features --- implemented}

The feature-engineering pipeline uses only current and past information.
The temperature-only feature set includes:

\begin{itemize}[leftmargin=2em,itemsep=2pt]
  \item Current temperature
  \item Temperature lag 1, 2, 4, 8, 12
  \item Temperature change
  \item Rolling mean, minimum, maximum, standard deviation
  \item Temperature trend
  \item Hour, minute of day
\end{itemize}

\noindent The multivariate research dataset additionally includes:

\begin{itemize}[leftmargin=2em,itemsep=2pt]
  \item Humidity
  \item CO\textsubscript{2}
  \item Lighting
  \item Wind
  \item Solar measurements
  \item Pyranometer values
\end{itemize}

\noindent \textbf{No future values are used as input features.} All
target shifts (30/60/90-minute) are derived after the feature matrix is
built to keep the input strictly historical.

% ====================================================================
\section{External models --- trained}

The external benchmark trained regression and classification models on
\emph{both} datasets (temperature-only and multivariate), giving:

\begin{center}
\begin{tabularx}{0.9\textwidth}{L{6.5cm}L{7cm}}
\toprule
\textbf{Model family} & \textbf{Models trained} \\
\midrule
Regression     & Ridge regression, small decision tree, small random forest \\
Classification & Logistic regression, decision-tree classifier, random-forest classifier \\
\bottomrule
\end{tabularx}
\end{center}

\noindent A total of \textbf{24 model artifacts} were generated
successfully across the two datasets and the two task families.

% ====================================================================
\section{External evaluation reports --- created}

The following evidence files were generated:

\begin{center}
\begin{tabular}{ll}
\toprule
\textbf{Format} & \textbf{File} \\
\midrule
Markdown & \texttt{DATASET\_AUDIT.md} \\
Markdown & \texttt{EDA\_REPORT.md} \\
Markdown & \texttt{EXTERNAL\_BENCHMARK\_REPORT.md} \\
JSON     & \texttt{dataset\_audit.json} \\
JSON     & \texttt{limitations.json} \\
CSV      & \texttt{baseline\_results.csv} \\
CSV      & \texttt{regression\_model\_comparison.csv} \\
CSV      & \texttt{classification\_model\_comparison.csv} \\
CSV      & \texttt{group\_validation\_report.csv} \\
\bottomrule
\end{tabular}
\end{center}

\noindent Plots were also generated (\texttt{evidence/external\_t15/plots/}) for:
\begin{itemize}[leftmargin=2em,itemsep=2pt]
  \item Dataset timeline
  \item Split row counts
  \item Regression MAE comparison
  \item Classification macro-F1 comparison
\end{itemize}

% ====================================================================
\section{Reproducible external workflow --- created}

The following PowerShell scripts were added:

\begin{center}
\begin{tabular}{ll}
\toprule
\textbf{Script} & \textbf{Step} \\
\midrule
\texttt{run\_external\_t15\_audit.ps1}        & Dataset audit \\
\texttt{build\_external\_t15\_dataset.ps1}    & Model-ready dataset build \\
\texttt{train\_external\_t15\_models.ps1}     & Model training \\
\texttt{evaluate\_external\_t15\_models.ps1}  & Model evaluation \\
\texttt{verify\_external\_t15\_phase.ps1}     & Phase verification \\
\bottomrule
\end{tabular}
\end{center}

\noindent The correct execution sequence is:
\begin{enumerate}[leftmargin=2em,itemsep=2pt]
  \item Raw data conversion (\texttt{tools/convert\_t15\_to\_thermal\_nexus.py})
  \item Dataset audit
  \item Model-ready dataset build
  \item Model training
  \item Model evaluation
  \item Phase verification
\end{enumerate}

% ====================================================================
\section{Git collaboration workflow --- prepared}

A clean Git branch was created: \textbf{\texttt{provat/external-t15-clean}}.
The branch was pushed successfully to GitHub and subsequently merged into
\texttt{master} (merge commit \texttt{089e195}; preview branch
\texttt{provat/external-t15-finalization} remains for review).

\noindent Large or private generated files were excluded through
\texttt{.gitignore}, including:
\begin{itemize}[leftmargin=2em,itemsep=2pt]
  \item Raw external data
  \item Processed external data
  \item Rebuilt data
  \item Generated model-ready CSV files
  \item Model \texttt{.joblib} files
  \item Temporary ZIP files
  \item Temporary build logs
\end{itemize}

\noindent This keeps the Pull Request smaller and prevents unnecessary or
unlicensed data from being published.

% ====================================================================
\section{Clean-clone testing --- performed}

A fresh temporary worktree was created to simulate what a teammate or
CI system receives after cloning the repository. The initial fresh-clone
result was:

\begin{center}
\begin{tabular}{lr}
\toprule
\textbf{Result} & \textbf{Value} \\
\midrule
Tests passed           & 47 \\
Tests failed           & 7 \\
Ruff                   & passed \\
Black                  & passed \\
\bottomrule
\end{tabular}
\end{center}

\noindent The seven failures were not model failures. They occurred because
some tests depended on local generated files such as:

\begin{itemize}[leftmargin=2em,itemsep=2pt]
  \item \texttt{ml/data/splits/validation.csv}
  \item \texttt{ml/data/external/t15/processed/...}
  \item \texttt{ml/data/external/t15/model\_ready/...}
\end{itemize}

\noindent These files were intentionally ignored by Git and therefore
absent from a clean clone.

% ====================================================================
\section{Test-isolation fix --- implemented}

The testing system is being changed so tests no longer depend on local
generated datasets. The changes include:

\begin{itemize}[leftmargin=2em,itemsep=2pt]
  \item Small test fixtures under \texttt{tests/fixtures/}
  \item Temporary test directories via \texttt{tmp\_path}
  \item Path replacement via \texttt{monkeypatch}
  \item Optional input/output paths in production functions
  \item Removal of hardcoded validation-file dependencies
  \item External T15 tests using small synthetic fixtures
  \item Optional integration markers for full external-data tests
  \item Protection against tests modifying repository data
\end{itemize}

\noindent Files modified for this work include:

\begin{center}
\begin{tabular}{ll}
\toprule
\textbf{Area} & \textbf{Files} \\
\midrule
Embedded deployment  & \texttt{estimate\_resources.py}, \texttt{generate\_golden\_vectors.py}, \\
                     & \texttt{prepare\_export.py} \\
Embedded tests       & \texttt{tests/test\_parity.py} \\
External T15         & \texttt{ml/external\_t15/dataset.py} \\
ML preprocessing     & \texttt{ml/preprocessing/split\_dataset.py} \\
Project config       & \texttt{pyproject.toml} \\
Test suites          & \texttt{tests/test\_dashboard\_phase.py}, \\
                     & \texttt{tests/test\_external\_t15\_pipeline.py}, \\
                     & \texttt{tests/test\_ml\_preprocessing.py}, \\
                     & \texttt{tests/test\_runtime\_simulation.py}, \\
                     & \texttt{tests/test\_training\_pipeline.py} \\
Fixtures             & \texttt{tests/fixtures/} \\
\bottomrule
\end{tabular}
\end{center}

% ====================================================================
\section{Current project status}

\begin{center}
\renewcommand{\arraystretch}{1.25}
\begin{tabularx}{0.92\textwidth}{L{9.5cm}L{4cm}}
\toprule
\textbf{Area} & \textbf{Status} \\
\midrule
Core simulator and reader pipeline          & Complete \\
SQLite data storage                        & Complete \\
Offline dashboard                          & Complete \\
KPI and energy analysis                    & Complete \\
Embedded export workflow                   & Complete \\
External T15 conversion                    & Complete \\
External dataset audit                     & Complete \\
Feature engineering                        & Complete \\
30/60/90-minute targets                    & Complete \\
Model training                             & Complete \\
Model evaluation                           & Complete \\
External reports                           & Complete \\
Scientific limitations                     & Complete \\
Clean Git branch                           & Complete \\
Branch pushed to GitHub                    & Complete \\
Fresh-clone problem identified             & Complete \\
Test-isolation changes                     & Implemented \\
\bottomrule
\end{tabularx}
\end{center}

\vspace{1em}
\noindent\textit{Note.} The four remaining failures in the current full
suite (\texttt{tests/test\_external\_t15\_pipeline.py}) are the expected
``clean-clone'' failures: they reference the gitignored raw
\texttt{NEW-DATA-*.T15.txt} sources and are intended to be marked as
integration tests in the final test-isolation PR.

\end{document}